\documentclass[12pt,a4paper,oneside]{article}
\usepackage[brazil]{babel}
\usepackage[utf8]{inputenc}
\usepackage{olymp}

\setcounter{problem}{0}

\begin{document}

\begin{problem}{Inscrição duplicada}{duplicada}{2}

As inscrições para mestrado e doutorado em Ciência da Computação na UFV foram abertas nesta semana.
O curso é \textit{multicampi}, podendo ser feito em Viçosa ou em Florestal.
No ato da inscrição, os candidatos escolhem o \textit{campus} que desejam se vincular.

Alguns candidatos acabam se inscrevendo nos dois \textit{campi} dando um trabalho duplicado para a comissão de seleção.

Faça um programa para descobrir quantos candidatos se inscreveram nos dois \textit{campi}.

%\Notes
%
%Observações, se tiver.

\InputFile

A entrada contém três linhas:
a primeira contém dois números inteiros, $N$ e $M$, representando a quantidade de inscrições em Viçosa e em Florestal respectivamente;
a segunda linha contém $N$ números inteiros $V_i$, os ID dos candidatos inscritos em Viçosa;
e a terceira linha contém $M$ números inteiros $F_i$, os ID dos candidatos inscritos em Florestal.

Restrições: $1 \leq N, M \leq 200$ e $1 \leq V, F \leq 1000000$,
e não há inscrições repetidos em um mesmo \textit{campus}.

\OutputFile

Seu programa deve gerar apenas uma linha de saída, contendo a quantidade de ID's que aparecem simultaneamente entre os inscritos em Viçosa e em Florestal.

%\Notes

\Example

\begin{example}
\exmp{
5 5
4 5 3 2 1
9 4 7 1 8
}{
2
}%
\end{example}

\begin{example}
\exmp{
5 5
89 26 37 42 11
26 11 89 42 37
}{
5
}%
\end{example}

\begin{example}
\exmp{
3 5
44 22 11
20 60 85 10 2
}{
0
}%
\end{example}

%Comentários sobre o exemplo, se necessário.
\small
{\textit {No primeiro exemplo, os candidatos 4 e 1 aparecem nas duas listas;\\
no segundo exemplo, todos os 5 aparecem nas duas listas; e, no último exemplo, nenhum.}}

\end{problem}

\end{document}